\documentclass[11pt]{article}
\usepackage{amsmath,amssymb,amsthm}
\usepackage{mathpartir}
\usepackage[margin=2.5cm]{geometry}
\usepackage{stmaryrd}

% ---- syntax -------------------------------------------------------------
\newcommand{\Ctx}{\mathsf{Ctx}}
\newcommand{\wf}[1]{#1\;\mathsf{ctx}}
\newcommand{\hastype}[3]{#1 \vdash #2 : #3}
\newcommand{\conveq}[4]{#1 \vdash #2 = #3 : #4}
\newcommand{\Pitype}[2]{\Pi(#1)\,#2}
\newcommand{\Lamterm}[2]{\lambda(#1)\,#2}
\newcommand{\App}[2]{#1\;#2}
\newcommand{\substone}[2]{#1[#2]}
\newcommand{\wkn}[1]{#1^{+}}
\newcommand{\U}{\mathsf{U}}
\newcommand{\Red}{\mathrel{\leadsto^{*}}}

% ---- domain model -------------------------------------------------------
\newcommand{\Bot}{\bot}
\newcommand{\PiC}[2]{\Pi(#1,#2)}        % Pi-type code
\newcommand{\FunEl}[1]{\lambda #1}      % function value with finite graph
\newcommand{\ev}[2]{#1\!\cdot\!#2}      % evaluation Eval(f,u)
\newcommand{\app}{\mathbin{+\!\!+}}     % append of finite functions
\newcommand{\comp}{\mathbin{\uparrow}}  % compatibility (common upper bound)
\newcommand{\Code}{\mathsf{Code}}
\newcommand{\FunCode}{\mathsf{FunCode}}
\newcommand{\rk}{\mathrm{rk}}
\newcommand{\Coh}{\mathsf{Coh}}
\newcommand{\Nat}{\mathbb{N}}
% stage-indexed relations
\newcommand{\lest}[1]{\mathrel{\leq_{#1}}}      % u <=_n v
\newcommand{\memst}[3]{#1 \mathrel{:_{#2}} #3}   % u :_n a   (\memst{u}{n}{a})
% validity
\newcommand{\Val}{\mathsf{Val}}
\newcommand{\TVal}{\mathsf{TVal}}
\newcommand{\EqVal}{\mathsf{EqVal}}
\newcommand{\EqTVal}{\mathsf{EqTVal}}

\theoremstyle{plain}
\newtheorem{theorem}{Theorem}
\newtheorem{lemma}{Lemma}
\newtheorem{corollary}{Corollary}
\theoremstyle{definition}
\newtheorem{definition}{Definition}

\title{A Stratified Domain Model for $\U:\U$ with $\Pi$-types:\\
       Typing, Validity, Adequacy, and $\Pi$-Injectivity}
\author{}
\date{}

\begin{document}
\maketitle

This note describes the $\Pi+\U$ fragment of the domain-theoretic model of
Coquand--Huber, in the form mechanised in the \texttt{MIN/} development.
It is a simplification of the $\Pi+\Sigma+\U$ presentation
(\texttt{sigma-rules.tex}): there are no $\Sigma$-types and no $\mathsf{Prop}$
universe.  The novelty here is methodological.  Rather than defining the
information order $\leq$ and the semantic membership $u : a$ directly---which
forces a \texttt{TERMINATING} pragma in Agda---we define them
\emph{by stratification}: a family of relations $\lest{n}$, $\memst{\cdot}{n}{\cdot}$
indexed by a stage $n\in\Nat$, each obtained from the previous one, together
with a \emph{stability} lemma that lets us pass to the limit.  The same idea
is used for the logical relation.  Every definition is then structurally
recursive on the stage index, so the whole development is free of
\texttt{TERMINATING} and \texttt{NO\_POSITIVITY} pragmas, of postulates, and
of holes.

\section{Judgments}

We use the two judgment forms (well-formedness of contexts is folded in):
\[
  \wf{\Gamma}\,, \qquad
  \hastype{\Gamma}{M}{A}\,, \qquad
  \conveq{\Gamma}{M}{N}{A}\,.
\]
Contexts are snoc-lists $\Gamma ::= \cdot \mid \Gamma, x{:}A$.  We write
$M\sigma$ for substitution and $\substone{B}{a}$ for the substitution of $a$
for the last variable.  We work in a type-in-type setting, $\U:\U$; the system
is therefore logically inconsistent (Girard's paradox), so we do not address
normalisation.  Adequacy, $\Pi$-injectivity, and subject reduction are
statements about the logical relation and are independent of consistency.

\section{Typing rules}

\begin{mathpar}
  \inferrule[\textsc{Var}]
    {\wf{\Gamma} \\ (x{:}A) \in \Gamma}
    {\hastype{\Gamma}{x}{A}}
  \and
  \inferrule[\textsc{U}]
    {\wf{\Gamma}}
    {\hastype{\Gamma}{\U}{\U}}
  \and
  \inferrule[\textsc{Conv}]
    {\hastype{\Gamma}{M}{A} \\
     \conveq{\Gamma}{A}{B}{\U} \\
     \hastype{\Gamma}{B}{\U}}
    {\hastype{\Gamma}{M}{B}}
  \\
  \inferrule[\textsc{Pi}]
    {\hastype{\Gamma}{A}{\U} \\
     \hastype{\Gamma, x{:}A}{B}{\U}}
    {\hastype{\Gamma}{\Pitype{x{:}A}{B}}{\U}}
  \and
  \inferrule[\textsc{Lam}]
    {\hastype{\Gamma}{A}{\U} \\
     \hastype{\Gamma, x{:}A}{B}{\U} \\
     \hastype{\Gamma, x{:}A}{M}{B}}
    {\hastype{\Gamma}{\Lamterm{x{:}A}{M}}{\Pitype{x{:}A}{B}}}
  \and
  \inferrule[\textsc{App}]
    {\hastype{\Gamma}{A}{\U} \\
     \hastype{\Gamma, x{:}A}{B}{\U} \\
     \hastype{\Gamma}{f}{\Pitype{x{:}A}{B}} \\
     \hastype{\Gamma}{a}{A}}
    {\hastype{\Gamma}{\App{f}{a}}{\substone{B}{a}}}
\end{mathpar}
The context-formation rule is $\wf{\cdot}$ and
$\hastype{\Gamma}{A}{\U} \Rightarrow \wf{(\Gamma, x{:}A)}$.

\section{Conversion rules}

\begin{mathpar}
  \inferrule[\textsc{Refl}]{\hastype{\Gamma}{M}{A}}{\conveq{\Gamma}{M}{M}{A}}
  \and
  \inferrule[\textsc{Sym}]{\conveq{\Gamma}{M}{N}{A}}{\conveq{\Gamma}{N}{M}{A}}
  \and
  \inferrule[\textsc{Trans}]
    {\conveq{\Gamma}{M}{N}{A} \\ \conveq{\Gamma}{N}{P}{A}}
    {\conveq{\Gamma}{M}{P}{A}}
  \and
  \inferrule[\textsc{Conv-Eq}]
    {\conveq{\Gamma}{M}{N}{A} \\ \conveq{\Gamma}{A}{B}{\U} \\ \hastype{\Gamma}{B}{\U}}
    {\conveq{\Gamma}{M}{N}{B}}
  \\
  \inferrule[\textsc{Beta}]
    {\hastype{\Gamma}{A}{\U} \\
     \hastype{\Gamma, x{:}A}{B}{\U} \\
     \hastype{\Gamma, x{:}A}{M}{B} \\
     \hastype{\Gamma}{a}{A}}
    {\conveq{\Gamma}{\App{(\Lamterm{x{:}A}{M})}{a}}{\substone{M}{a}}{\substone{B}{a}}}
  \and
  \inferrule[\textsc{Pi-Cong}]
    {\conveq{\Gamma}{A}{A'}{\U} \\ \conveq{\Gamma, x{:}A}{B}{B'}{\U}}
    {\conveq{\Gamma}{\Pitype{x{:}A}{B}}{\Pitype{x{:}A'}{B'}}{\U}}
  \and
  \inferrule[\textsc{App-Fun}]
    {\conveq{\Gamma}{f}{f'}{\Pitype{x{:}A}{B}} \\ \hastype{\Gamma}{a}{A} \\ \cdots}
    {\conveq{\Gamma}{\App{f}{a}}{\App{f'}{a}}{\substone{B}{a}}}
  \and
  \inferrule[\textsc{App-Arg}]
    {\hastype{\Gamma}{f}{\Pitype{x{:}A}{B}} \\ \conveq{\Gamma}{a}{a'}{A} \\ \cdots}
    {\conveq{\Gamma}{\App{f}{a}}{\App{f}{a'}}{\substone{B}{a}}}
  \and
  \inferrule[\textsc{Funext}]
    {\conveq{\Gamma, x{:}A}{\App{\wkn{f}}{x}}{\App{\wkn{g}}{x}}{B} \\
     \hastype{\Gamma}{f}{\Pitype{x{:}A}{B}} \\
     \hastype{\Gamma}{g}{\Pitype{x{:}A}{B}} \\ \cdots}
    {\conveq{\Gamma}{f}{g}{\Pitype{x{:}A}{B}}}
\end{mathpar}
\textsc{Funext} is the only extensionality principle; it is validated by the
model, where two functions with the same finite graph are identified.

\section{Head reduction}

Single-step head reduction $M \leadsto_1 N$ is $\beta$ at the head, closed under
the application head:
\begin{mathpar}
  \inferrule[\textsc{HR-Beta}]{ }
    {\App{(\Lamterm{x{:}A}{M})}{a} \leadsto_1 \substone{M}{a}}
  \and
  \inferrule[\textsc{HR-App}]{M \leadsto_1 M'}{\App{M}{a} \leadsto_1 \App{M'}{a}}
\end{mathpar}
We write $M \Red N$ for its reflexive--transitive closure.  The model tracks
only head normal forms: type validity stores $T \Red \Pitype{x{:}A}{B}$, and
adequacy reconstructs convertibility from these reductions.  Reduction under
binders is never used.

\section{The domain model}

\subsection{Finite elements}

Codes (finite elements of the Scott domain $D$) and finite functions are given
inductively:
\[
  u,v,a \in \Code ::= \Bot \mid \U \mid \FunEl g \mid \PiC{a}{f},
  \qquad
  f,g \in \FunCode ::= [\,] \mid (u,v) :: f.
\]
$\Bot$ carries no information, $\U$ is the universe code, $\FunEl g$ is a
function value with finite graph $g$ (a finite list of input/output pairs),
and $\PiC{a}{f}$ is a $\Pi$-type code with domain code $a$ and codomain
function $f$.  Each $u$ has a \emph{rank} $\rk(u)\in\Nat$
(with $\rk(\PiC{a}{f}),\rk(\FunEl g) = 1+\max(\cdots)$), measuring how many
iterative stages are needed to build it.

\subsection{Supremum and evaluation}

The binary supremum combines elements with the same head constructor and is
$\Bot$ otherwise:
\[
  \Bot \sqcup u = u,\quad \U\sqcup\U = \U,\quad
  \FunEl g \sqcup \FunEl h = \FunEl{(g\app h)},\quad
  \PiC{a}{f}\sqcup\PiC{b}{g} = \PiC{a\sqcup b}{f\app g}.
\]
We write $u \comp v$ (``$u$ and $v$ are \emph{compatible}'') when they have a
common upper bound; on coherent elements this is exactly ``same head, and
recursively compatible''.  Evaluation of a finite function takes the supremum
of the outputs whose keys are below the argument:
\[
  \ev f u \;=\; \bigsqcup\{\, v_i \mid (u_i,v_i)\in f,\ u_i \leq u \,\}.
\]

\subsection{Stages: a stratified construction}
\label{sec:stages}

The domain $D$ is the limit (ideal completion) of an increasing chain
$D_0 \subseteq D_1 \subseteq \cdots$ of finite approximations.  A central
methodological point is that \emph{we do not reason through projections}
$p_n : D \to D_n$: the naive transfer ``$u : a \Rightarrow p_n u : p_n a$''
\emph{fails}.  Instead, every finite element of $D$ is presented as the union
of an increasing sequence of finite elements of the $D_n$, and we work
entirely with these finite elements.  At each stage $D_n$ is a
\emph{conditional sup-semilattice} (compatible elements have a least upper
bound).

Concretely we define, by induction on the stage $n$, a relation $\lest{n}$ on
codes and a relation $\memst{\cdot}{n}{\cdot}$ (``$u :_n a$''), each defined
``extensionally'' from the level below.  The level-$(n{+}1)$ clauses mention
only level-$n$ data, so the definitions are structurally recursive in $n$.
For example, type membership at a $\Pi$-code is
\begin{equation}
  \label{eq:stage-pi}
  \memst{\PiC{a}{f}}{n+1}{\U}
  \;\;\Longleftrightarrow\;\;
  \memst{a}{n}{\U}
  \;\wedge\;
  \bigl(\forall x\,\exists y\leq x.\ \memst{y}{n}{a}\ \wedge\ \ev f x = \ev f y\bigr)
  \;\wedge\;
  \bigl(\forall x.\ \memst{(\ev f x)}{n}{\U}\bigr),
\end{equation}
expressing that $\PiC{a}{f}$ is a well-formed type at stage $n{+}1$ iff its
domain is a type, its graph is determined by finitely many points each lying
in $a$, and all its outputs are types --- all checked one stage down.  The
order on functions is likewise extensional, $f \lest{n} g \iff \forall x.\,
\ev f x \lest{n} \ev g x$.

\begin{lemma}[Stability]\label{lem:stab}
  For codes $u,a$ with ranks bounded by $n$,
  \[
    u \lest{n} v \iff u \lest{n+1} v
    \qquad\text{and}\qquad
    \memst{u}{n}{a} \iff \memst{u}{n+1}{a}.
  \]
\end{lemma}
Stability is proved by induction on $n$, using that a function value has a
finite representation $(u_1,v_1),\dots,(u_m,v_m)$ with $\memst{u_i}{n}{a}$ and
$\memst{v_i}{n}{\U}$.  By stability the stage stops mattering above the rank,
and we define the limit relations
\[
  u \leq v \;:=\; u \lest{n} v
  \qquad\text{and}\qquad
  u : a \;:=\; \memst{u}{n}{a}
  \qquad\text{for any } n > \max(\rk u, \rk v),\ n > \max(\rk u, \rk a),
\]
which are well-defined by Lemma~\ref{lem:stab}.  Because the construction is
recursion on $n$, the Agda kernel needs no \texttt{TERMINATING} pragma.

\subsection{The information order $\leq$}

Unfolding the limit, $\leq$ is the expected structural order: $\Bot \leq u$ for
all $u$; $\U \leq \U$; on compound codes
\[
  \FunEl g \leq \FunEl h \iff
  \forall (u_i,v_i)\in g.\ v_i \leq \ev h {u_i},
  \qquad
  \PiC{a}{f} \leq \PiC{b}{g} \iff a\leq b \ \wedge\ f \leq g,
\]
and codes with different heads are incomparable.

\subsection{Coherence}

A finite function is \emph{coherent} when compatible keys force compatible
values: $\Coh(g)$ requires that each key and value be coherent and non-$\Bot$,
and that $u_i \comp u_j \Rightarrow v_i \comp v_j$.  On codes,
$\Coh(\PiC{a}{f}) = \Coh(a)\wedge\Coh(f)$, while $\Bot$ and $\U$ are always
coherent.  Coherence makes evaluation monotone:
$u \leq v$ and $\Coh(g)$ imply $\ev g u \leq \ev g v$.

\subsection{The membership relation $u : a$}

The relation $u : a$ (``the finite element $u$ approximates an element of the
type $a$'') unfolds, on heads, to:
\[
  \begin{aligned}
    \Bot : a &\iff a : \U
      &&\text{(}\Bot\text{ approximates every type),}\\
    \U : \U &\iff \top, \\
    \PiC{a}{f} : \U &\iff a : \U \ \wedge\ (\forall (u_i,v_i)\in f.\ u_i : a \wedge v_i : \U)\ \wedge\ \Coh(f),\\
    \FunEl g : \PiC{a}{f} &\iff \bigl(\forall (u_i,v_i)\in g.\ u_i : a \wedge v_i : \ev f {u_i}\bigr)
      \ \wedge\ \Coh(g)\ \wedge\ \PiC{a}{f} : \U,
  \end{aligned}
\]
all other head combinations being trivial.  The last clause is the meaning of
``\,$\FunEl g$ is a function in $\PiC{a}{f}$\,'': every entry of its graph sends
a domain point $u_i : a$ to a value $v_i$ of the corresponding fibre
$\ev f{u_i}$, the graph is coherent, and $\PiC{a}{f}$ is itself a type.

\subsection{Main properties}

All proved in Agda, pragma-free and postulate-free:
\begin{itemize}
  \item \textbf{Order.} $\leq$ is reflexive on coherent codes, transitive, and
    $\Bot$ is least.  Compatible codes have a least upper bound: if $u\comp v$
    then $u\leq u\sqcup v$, $v\leq u\sqcup v$, and $u\sqcup v$ is below any
    common bound.  \textbf{(Comp-down)} $u\leq u'$ and $u'\comp v$ imply
    $u\comp v$.  Evaluation is monotone in both function and argument.
  \item \textbf{Type-hood.} $\Bot : a \iff a : \U$, i.e.\ $a$ is a type.
  \item \textbf{Compatible supremum.} If $u_1 : a$ and $u_2 : a$ and
    $u_1 \comp u_2$, then $u_1 \sqcup u_2 : a$.  (Membership is closed under
    conditional sups; this is the semilattice property at the membership
    level.)
  \item \textbf{Monotonicity.} If $u : a$ and $a \leq b$ and $b : \U$, then
    $u : b$.  (A member of a smaller type is a member of a larger one.)
  \item \textbf{Evaluation in $\U$.} If $f$'s values lie in $\U$ over $a$ and
    $x : a$, then $\ev f x : \U$.
\end{itemize}

\section{The logical relation (validity)}

The logical relation connects the model to the syntax.  It is built from four
mutually-defined predicates, again \emph{stratified by a stage} $n$ (the
\texttt{ValidityStratified}/\texttt{GoodStage} construction) and collapsed by a
stability lemma, so that it too is structurally recursive on $n$.

\subsection{Term and type validity}

\begin{itemize}
  \item $\Val\;\Gamma\;M\;A\;u\;a$: the term $M:A$ in $\Gamma$ is valid at value
    code $u$ and type code $a$.
  \item $\TVal\;\Gamma\;M\;a$: the type $M$ is valid at code $a$.
  \item $\EqVal$, $\EqTVal$: the corresponding equality (two-sided) relations.
\end{itemize}
The defining clauses mirror the membership relation.  At the universe,
$\Val\;\Gamma\;M\;A\;u\;\U$ (for $u\in\{\U,\PiC{b}{f}\}$) asks for
$\TVal\;\Gamma\;A\;\U$ and $\TVal\;\Gamma\;M\;u$.  At a $\Pi$-code,
$\TVal\;\Gamma\;M\;\PiC{b}{f}$ unpacks to a head reduction
$M \Red \Pitype{x{:}A}{B}$ with $\hastype{\Gamma}{A}{\U}$,
$\hastype{\Gamma,x{:}A}{B}{\U}$, domain validity $\TVal\;\Gamma\;A\;b$, and an
\emph{edge function}: for every selection $(u,v)$ from $f$ and well-typed
argument $N$ with $\Val\;\Gamma\;N\;A\;u\;b$, one has
$\TVal\;\Gamma\;\substone{B}{N}\;v$.  Dually, $\Val$ at $(\FunEl g,\PiC{b}{f})$
asks, for every selection $(u,v)$ from $g$ and such $N$, that
$\Val\;\Gamma\;(\App M N)\;\substone{B}{N}\;v\;(\ev f u)$.  A \emph{selection}
of $f$ takes or skips each entry, combining the taken pairs by $\sqcup$, and
yields $(u,v)$ with $u \leq \ev f u$ built from subterms of $f$.

\subsection{Stages and stability for validity}

As in Section~\ref{sec:stages}, $\Val,\TVal,\EqVal,\EqTVal$ are defined by
induction on $n$.  One proves, by induction on $n$, that if the four relations
satisfy all the required properties at stage $n$ then they satisfy them at
stage $n{+}1$; among these properties is \emph{monotonicity of type validity}:
if $\TVal\;\Gamma\;A\;a$ and $a' \leq a$ with $a' : \U$, then
$\TVal\;\Gamma\;A\;a'$.  Monotonicity then drives stability: for $a$ at stage
$n$, $\TVal_n\;\Gamma\;A\;a \iff \TVal_{n+1}\;\Gamma\;A\;a$, and similarly for
the other three.  The collapsed relations are the stage limits.  (The edge
functions store $\Val$ in negative position; defining the records in a module
parameterised by the four relations keeps the positivity checker satisfied
with no \texttt{NO\_POSITIVITY} pragma.)

\subsection{Main properties}

\begin{itemize}
  \item \textbf{Monotonicity (down/up).} Validity transports along $\leq$ on the
    type code, in both directions where the codes are types.
  \item \textbf{Supremum.} Validity is closed under compatible sups of codes,
    matching the membership-level property.
  \item \textbf{Reflexivity / symmetry / transitivity} of the equality
    relations $\EqVal,\EqTVal$, and the bundling that exposes the syntactic
    \textsc{HasType}/\textsc{ConvTm} witnesses at the leaves.
  \item \textbf{Head expansion.} Validity is preserved by head reduction of the
    subject (used in the $\beta$ case of adequacy).
\end{itemize}

\section{Adequacy}

\begin{theorem}[Adequacy]\label{thm:adequacy}
  If $\hastype{\Gamma}{M}{A}$ then, for any coherent environment $\rho$ with a
  valid substitution $\sigma$ into a context $H$, and codes $u,a$ such that $M$
  evaluates to $u$ at $\rho$, $A$ evaluates to $a$ at $\rho$, and $u : a$, we
  have $\Val\;H\;\sigma(M)\;\sigma(A)\;u\;a$.
\end{theorem}
The proof is by induction on the typing derivation, with a parallel statement
for conversion ($\conveq{\Gamma}{M}{N}{A}$ implies $\EqVal$) and a
two-substitution variant ($M$ under two related substitutions yields $\EqVal$).
Injectivity and subject reduction are corollaries, obtained by evaluating at
the trivial bottom environment $\rho_\bot$ (under which $\Pi B\,F$ evaluates to
$\PiC{\Bot}{[\,]}$) and reading the head-reduction and conversion data off the
logical relation.

\section{$\Pi$-injectivity}

\begin{theorem}[$\Pi$-injectivity]\label{thm:pi-inj}
  If $\conveq{\Gamma}{A_0}{\Pitype{x{:}B_1}{F_1}}{\U}$, then there exist
  $B_0, F_0$ with
  \begin{enumerate}
    \item $A_0 \Red \Pitype{x{:}B_0}{F_0}$,
    \item $\conveq{\Gamma}{B_0}{B_1}{\U}$,
    \item $\conveq{\Gamma, x{:}B_0}{F_0}{F_1}{\U}$.
  \end{enumerate}
\end{theorem}

\begin{corollary}[$\Pi$--$\Pi$ injectivity]
  If $\conveq{\Gamma}{\Pitype{x{:}A_0}{B_0}}{\Pitype{x{:}A_1}{B_1}}{\U}$ then
  $\conveq{\Gamma}{A_0}{A_1}{\U}$ and
  $\conveq{\Gamma, x{:}A_0}{B_0}{B_1}{\U}$.
\end{corollary}

\begin{proof}[Proof outline]
  By presupposition $\hastype{\Gamma}{A_0}{\U}$.  Conversion soundness at
  $\rho_\bot$ transfers the evaluation of $\Pitype{x{:}B_1}{F_1}$ to
  $\PiC{\Bot}{[\,]}$ onto $A_0$.  Bundled equality adequacy yields $\EqVal$ at
  the code $\PiC{\Bot}{[\,]}$, whose type-validity component records the head
  reduction $A_0 \Red \Pitype{x{:}B_0}{F_0}$ together with the domain and
  codomain conversions.  Since $\Pi$ is a head normal form, reduction
  uniqueness gives the stated equalities.
\end{proof}

\medskip\noindent\textbf{Subject reduction} (a corollary, in
\texttt{MIN/SubjectReduction.agda}): if $\hastype{\Gamma}{M}{A}$ and
$M \leadsto_1 N$ then $\hastype{\Gamma}{N}{A}$.  The only nontrivial case is
$\beta$, where $\Pi$-injectivity inverts the typing of the redex's
$\lambda$-abstraction through the conversion.

\end{document}
